\documentclass[11pt]{article}

\usepackage[margin=1in]{geometry}
\usepackage{amsmath,amssymb,bm}
\usepackage{booktabs}
\usepackage{hyperref}

\title{Action-kernel behavioural model\\
\large\textit{fit\_action\_kernel} / IBL \texttt{behavior\_models.ActionKernel}}
\author{prior-mech-analysis notes}
\date{2026-07-22}

\begin{document}
\maketitle

\section{Context}

For act-conditioned analyses we currently estimate priors with a
\textbf{fixed} learning rate \(\alpha = 0.2\) via
\texttt{action\_kernel\_priors}, applied the same way on every session
before pooling into the supersession.

\texttt{fit\_action\_kernel} (in
\texttt{scripts/simulate\_synthetic\_choices.py}) instead fits the full
IBL \texttt{ActionKernel} choice model by MCMC on one session's real
behaviour and returns a posterior-mean parameter vector. Open question
(journal 2026-07-20): whether those session-specific fits should replace
the fixed-\(\alpha\) prior for act labels before supersession pooling.

\paragraph{Data.} Fitting uses
\[
(\text{actions},\ \text{stimuli},\ \text{stim\_side})
= \texttt{format\_data}(\texttt{trials}),
\]
i.e.\ choices \(a_t\in\{-1,+1\}\) (0 = no-go), signed contrast
\(c_t\in[-1,1]\), and stimulus side \(s_t\in\{-1,+1\}\).
Task \(\texttt{probabilityLeft}\) is \emph{not} a covariate in this
likelihood: the prior is the evolving action kernel.

\section{Parameters (\texttt{single\_zeta=True})}

\begin{center}
\begin{tabular}{@{}llll@{}}
\toprule
Symbol & Code name & Role & Bounds \\
\midrule
\(\alpha\) & \texttt{alpha} & Learning rate of the action-history prior & \([0,1]\) \\
\(\zeta\) & \texttt{zeta} & Sensory noise (shared across sides) & \([0,1]\) \\
\(\ell_+\) & \texttt{lapse\_pos} & Lapse when stim side is left (\(s_t>0\)) & \([0,0.5]\) \\
\(\ell_-\) & \texttt{lapse\_neg} & Lapse when stim side is right (\(s_t<0\)) & \([0,0.5]\) \\
\bottomrule
\end{tabular}
\end{center}

Posterior mean
\(\hat{\bm\theta}=(\hat\alpha,\hat\zeta,\hat\ell_+,\hat\ell_-)\)
is what \texttt{get\_parameters(parameter\_type="posterior\_mean")}
returns after MCMC.

\section{Model equations}

\paragraph{Notation (probabilities).}
Three different probabilities appear below---do not conflate them:
\begin{itemize}
\item \(\pi_t^{\mathrm{L}}\) (and \(p_t\) in the analysis helper):
  \textbf{prior belief} that the next / current trial favours left,
  from the action-kernel EMA over past choices.
  \emph{Not} the instantaneous choice probability.
\item \(\tilde P(\mathrm{L}_t)\):
  \textbf{pre-lapse probability of choosing left}
  given current contrast \(c_t\) and prior \(\pi_t^{\mathrm{L}}\).
\item \(P(\mathrm{L}_t)\):
  \textbf{probability of choosing left after lapses}
  (what is scored in the likelihood and sampled in simulation).
  Then \(P(\mathrm{R}_t)=1-P(\mathrm{L}_t)\).
\end{itemize}

\subsection{Action-kernel prior}

Let \(\bm\pi_t=(\pi_t^{\mathrm{R}},\pi_t^{\mathrm{L}})\) be a probability
vector over \{right, left\}, initialised at
\(\bm\pi_0=(1/2,1/2)\). With learning rate \(\alpha\),

\begin{equation}
\label{eq:prior-update}
\bm\pi_t
=
\begin{cases}
(1-\alpha)\,\bm\pi_{t-1}
+ \alpha\,\mathbf{e}_{a_{t-1}}
& \text{if } a_{t-1}\neq 0,\\[0.4em]
\bm\pi_{t-1}
& \text{if } a_{t-1}=0\ \text{(no-go)},
\end{cases}
\end{equation}
where \(\mathbf{e}_{a}\) is the one-hot indicator of the previous choice
(\(a=+1\) \(\Rightarrow\) left, \(a=-1\) \(\Rightarrow\) right).
The scalar left-prior used below is \(\pi_t^{\mathrm{L}}=\bm\pi_t[1]\).

\paragraph{Range.}
The continuous left-prior is a probability:
\[
\pi_t^{\mathrm{L}}\in(0,1)
\quad\text{(equivalently }p_t\in(0,1)\text{ in \eqref{eq:simple-prior})}.
\]
It is initialised at \(1/2\) and stays in \((0,1)\) under the EMA update
(convex combination of a probability and \(\{0,1\}\)).
In \texttt{ActionKernel}, values are additionally clipped to
\([10^{-7},\,1-10^{-7}]\) before entering \eqref{eq:combine}
(so \(\Phi^{-1}(1-\pi)\) stays finite).
For act L/R \emph{splits} in analysis, \texttt{action\_kernel\_priors}
binarises the continuous prior to the discrete labels
\(\{0.8,\,0.2\}\) (\(\ge 1/2\mapsto 0.8\), else \(0.2\))---those are
conditioning labels, not the continuous \(\pi_t^{\mathrm{L}}\) used inside
the fitted choice model.

\paragraph{Analysis helper (fixed \(\alpha\) only).}
Our \texttt{action\_kernel\_priors} implements the same EMA on a scalar
\(p_t\):
\begin{equation}
\label{eq:simple-prior}
p_0 = \tfrac12,\qquad
p_t = \alpha\,\mathbf{1}[a_{t-1}>0] + (1-\alpha)\,p_{t-1},
\end{equation}
then binarises \(p_t\mapsto\{0.8,0.2\}\) for L/R-prior splits.
It does \emph{not} fit \(\zeta\) or lapses.

\subsection{Sensory likelihood combined with prior}

Signed contrast \(c_t\) is combined with the prior belief
\(\pi_t^{\mathrm{L}}\) under sensory noise \(\zeta\)
(\texttt{combine\_lkd\_prior}; fixed \(\sigma=0.49\))
to produce the \emph{choice} probability before lapses:
\begin{equation}
\label{eq:combine}
\begin{aligned}
\sigma^\star
&=
\bigl(\sigma^{-2}+\zeta^{-2}\bigr)^{-1/2},
\\
\tilde P(\mathrm{L}_t)
&\equiv
\tilde P(\text{choose left}\mid c_t,\pi_t^{\mathrm{L}},\zeta)
\\
&=
\Phi\!\left(
\frac{c_t}{\zeta}
-
\frac{\zeta}{\sigma^\star}\,
\Phi^{-1}\!\bigl(1-\pi_t^{\mathrm{L}}\bigr)
\right),
\end{aligned}
\end{equation}
where \(\Phi\) is the standard normal CDF.
(Equivalently
\(\tilde P(\mathrm{L}_t)
=\Phi\bigl(c_t/\zeta+(\zeta/\sigma^\star)\Phi^{-1}(\pi_t^{\mathrm{L}})\bigr)\)
since \(\Phi^{-1}(1-\pi)=-\Phi^{-1}(\pi)\).)

\subsection{Lapses}

Side-dependent lapse rate
\begin{equation}
\ell_t
=
\begin{cases}
\ell_+ & s_t > 0\ \text{(stim left)},\\
\ell_- & s_t < 0\ \text{(stim right)},\\
\tfrac12(\ell_++\ell_-) & s_t=0,
\end{cases}
\end{equation}
mixes the pre-lapse \emph{choice} probability with chance:
\begin{equation}
\label{eq:lapse}
P(\mathrm{L}_t)
\equiv
P(\text{choose left}\mid c_t,\pi_t,\bm\theta)
=
(1-\ell_t)\,\tilde P(\mathrm{L}_t)
+
\frac{\ell_t}{2}.
\end{equation}
With probability \(1-\ell_t\) the animal follows \(\tilde P\);
with probability \(\ell_t\) it guesses \(50/50\)
(hence the \(\ell_t/2\) term).
Always \(P(\mathrm{R}_t)=1-P(\mathrm{L}_t)\).

\subsection{Likelihood}

For observed choices \(\{a_t\}\),
\begin{equation}
\label{eq:lkd}
\mathcal{L}(\bm\theta)
=
\prod_{t:\,a_t\in\{\pm1\}}
P(a_t\mid c_t,\bm\pi_t,\bm\theta).
\end{equation}
MCMC targets the posterior of \(\bm\theta\) under this likelihood
(box priors on the bounds above). Early stopping uses
Gelman--Rubin diagnostics across chains.

\section{Simulation (\texttt{simulate\_choices})}

Given fixed \((c_t,s_t)\) and fitted \(\hat{\bm\theta}\),
\texttt{ActionKernel.simulate} runs the same forward pass
\eqref{eq:prior-update}--\eqref{eq:lapse}, but updates \(\bm\pi_t\) from
\textbf{previously simulated} choices (generative), and samples
\(a_t\sim\mathrm{Bernoulli}\bigl(P(\mathrm{L}_t)\bigr)\) each trial.
Stim/side (hence the empirical block schedule) stay fixed; only choices
are random.

\section{Relation to current analysis pipeline}

\begin{itemize}
\item \textbf{Act prior labels (current):}
  \eqref{eq:simple-prior} with \(\alpha=0.2\) on each session, then pool.
\item \textbf{ActKernel choice null (\texttt{--actkernel-choice-null}):}
  fit \(\bm\theta\) per \texttt{eid}, then BWM-style synthetic sessions
  (\texttt{generate\_pseudo\_session} stim/blocks + \texttt{simulate\_choices}
  under \(\hat{\bm\theta}\)); null labels are synthetic choices at the real
  session's stratified \texttt{elig\_idx}.
\item \textbf{Open:} replace fixed-\(\alpha\) labels with
  session-fitted \(\hat\alpha\) (or full kernel-derived priors) before
  supersession pooling.
\end{itemize}

\section*{References}

\begin{itemize}
\item Package:
  \url{https://github.com/int-brain-lab/behavior_models}
  (\texttt{ActionKernel}, \texttt{utils.combine\_lkd\_prior});
  vendored as submodule \texttt{third\_party/behavior\_models}.
\item Local wrapper:
  \texttt{scripts/simulate\_synthetic\_choices.py}
  (\texttt{fit\_action\_kernel}, \texttt{simulate\_choices},
  \texttt{synthetic\_sessions\_from\_trials}).
\item Analysis prior helper:
  \texttt{block\_analysis\_allsplits.action\_kernel\_priors}
  (\(\alpha=0.2\)).
\end{itemize}

\end{document}
